\documentclass[a4paper,fontset=none]{ctexart}

\usepackage{anyfontsize}
\usepackage{amsmath,amssymb,mathrsfs,wasysym}
\usepackage{autobreak}
\allowdisplaybreaks
\usepackage[T1]{fontenc}
\usepackage{textcomp}
\usepackage[libertine,vvarbb]{newtxmath}
\usepackage[scr=rsfso,frak=euler,bb=ams]{mathalfa}
\usepackage{bm}
\usepackage{indentfirst}
\setlength{\parindent}{0em}
\usepackage{parskip}
\usepackage{geometry}
\geometry{top=3cm,bottom=3cm,left=2cm,right=2cm}
\usepackage{fontspec}
\setmainfont{Libertinus Serif}
\setsansfont{Libertinus Sans}
\setmonofont{JetBrains Mono}
\setCJKmainfont{SimSun}[BoldFont=Noto Sans CJK SC Medium]

\begin{document}

    \title{\textbf{原子物理学作业}}
    \author{1900011604\ \ 张植竣\ \ 北京大学}
    \date{2021年3月8日}

    \maketitle

    \begin{itemize}

        \subsection*{思考题}

        \item[1.1]
        $1100\,\mathrm{K}$下，由Wien位移定律，辐射极大对应的波长在：
        \[
        \lambda_{\mathrm{max}} = \frac{b}{T} = 2.636\,\mu\mathrm{m}
        \]
        这对应于中红外波段，可见光部分呈现红色。

        由灰体辐射本领的谱密度：
        \[
        r(\nu,T) = B(\nu,T)\alpha(\nu)
        \]
        其中$B(\nu,T)$为黑体辐射本领的谱密度，$\alpha(\nu)$为吸收系数。

        可知：同一温度下，可见光波段黑体辐射本领的谱密度$B_\mathrm{v}(T)$相同，由于铁块的吸收系数远大于近乎透明的石英，其辐射本领也强，发出红色光辉，而石英在可见光波段的吸收系数$\alpha \approx 0$，辐射本领也弱，看上去不发光。

        \item[1.5]
        人类在从类人猿进化为人的过程中主要在阳光照耀的白昼活动，生物进化使眼睛最灵敏的波长与太阳光谱中辐射能最大的波长基本相符的个体最有可能存活，这是“适者生存”的需要。

        \item[1.10]
        Compton散射的前提是可以将散射原子中的电子看成是自由的和静止的。即光子的能量$E$远大于电子的逸出功$W$和动能$E_k$。电子的逸出功$W$可以用光电效应的红限波长估算，此时$\lambda_1 \approx 300\,\mathrm{nm}$，由
        \[
        E = \frac{hc}{\lambda}
        \]
        得$E \approx 4.13\,\mathrm{eV}$，而电子的动能小于其逸出功。Compton散射所用的电磁波是X射线，$\lambda_2 \approx 0.1\,\mathrm{nm}$，即$E \approx 1.24\times10^4\,\mathrm{eV}$，比逸出功约高了4个数量级，而可见光$\lambda_3 \approx 555\,\mathrm{nm}$，其对应的能量$E \approx 2.23\,\mathrm{eV}$，甚至小于部分材料的逸出功，自然不会发生Compton散射，可见光波长远大于散射光与入射光的波长差$\lambda_\mathrm{C} = 0.0024\,\mathrm{nm}$，也无法观测到Compton效应。综上，可见光不能用来做这个实验。

        \item[1.11]
        如果自由电子只吸收而不发射光子，将违背能量、动量守恒定律，故Compton散射中必然仍有光子散射出去。但是光电效应中电子还受到原子核的作用力，不是自由电子，吸收光子而不发射光子会增加电子动能但不违背能量、动量守恒定律。

        \subsection*{习题}

        \item[1.2]
        由题意得：
        \[
        \frac{4\pi R_{\astrosun}^2 \cdot \sigma T_{\astrosun}^4}{4\pi d^2} = 1.94\,\mathrm{cal/(cm^2 \cdot min)} = 1352.83\,\mathrm{W/m^2}
        \]
        其中$R_{\astrosun}$为太阳半径，$d$为日地距离，$T_{\astrosun}$为太阳表面的有效温度。

        又因为太阳的角半径：
        \[
        \theta = 0.00465\,\mathrm{rad} \approx \frac{R_{\astrosun}}{d}
        \]
        则有：
        \[
        \theta^2\sigma T_{\astrosun}^4 = 1352.83\,\mathrm{W/m^2}
        \]
        解得太阳表面的有效温度为：
        \[
        T_{\astrosun} = 5763.43\,\mathrm{K}
        \]

        \item[1.8(1)]
        由Einstein光电效应公式：
        \[
        \frac{hc}{\lambda} = W + E_\mathrm{k} = W + eU
        \]
        出射最快的光电子能量为：
        \[
        E_\mathrm{k} = \frac{hc}{\lambda} - W = 2.0\,\mathrm{eV}
        \]

        \item[1.8(2)]
        遏止电压为：
        \[
        U = \frac{E_\mathrm{k}}{e} = 2.0\,\mathrm{V}
        \]

        \item[1.8(3)]
        令$E_\mathrm{k} = 0$，得铝的截止波长为：
        \[
        \lambda_0 = \frac{hc}{E_\mathrm{k}} = 295.2\,\mathrm{nm}
        \]

        \item[1.8(4)]
        每个光子的平均能量为：
        \[
        E_0 = \frac{hc}{\lambda} = 9.932\times10^{-19}\,\mathrm{J}
        \]
        则单位时间打到单位面积上的平均光子数为：
        \[
        N = \frac{F\Delta t\Delta S}{E_0} = 2.014\times10^{18}
        \]

        \item[1.14]
        单个光子的动量大小为$p_0 = \dfrac{h\nu}{c}$，则对于全反射镜面，一个光子反冲传递给镜面的动量是$2p_0$。

        镜面上的辐射压强，即单位面积的受力为：
        \[
        P = \frac{1}{\Delta S}\cdot\frac{\mathrm{d}p}{\mathrm{d}t} = 2Np_0 = 2N\frac{h\nu}{c}
        \]

        \item[1.17(1)]
        首先计算Compton散射的规律，根据能量、动量守恒和相对论能动量关系可以得到下面四个方程：
        \[
        mc^2 + h\nu_0 = h\nu + E
        \]
        \[
        \frac{h\nu_0}{c} = \frac{h\nu}{c}\cos\alpha + p\cos\beta
        \]
        \[
        \frac{h\nu}{c}\sin\alpha = p\sin\beta
        \]
        \[
        E^2 = (pc)^2 + \left(mc^2\right)^2
        \]
        其中$\nu_0$和$\nu$分别是光子散射前后的频率，$\alpha$和$\beta$分别是散射后光子和电子相对于光子入射方向的散射角，$E$和$p$分别是电子的能量和动量。解得：
        \[
        \frac{1}{\nu} - \frac{1}{\nu_0} = \frac{h}{mc^2}(1-\cos\alpha)
        \]
        即：
        \[
        \Delta\lambda = \frac{c}{\nu} - \frac{c}{\nu_0} = \frac{h}{mc}(1-\cos\alpha)
        \]
        故对于$\lambda_0 = 0.100\,\mathrm{nm}$的X射线束，波长偏移为：
        \[
        \Delta\lambda = \frac{h}{mc}(1-\cos\alpha) = 0.002426\,\mathrm{nm}
        \]
        对于$\lambda_0 = 0.0188\,\mathrm{nm}$的$\gamma$射线束，波长偏移仍为：
        \[
        \Delta\lambda = \frac{h}{mc}(1-\cos\alpha) = 0.002426\,\mathrm{nm}
        \]

        \item[1.17(2)]
        注意到$\Delta\lambda$只与光子散射角$\alpha$有关，与入射光频率$\nu_0$无关，而$\nu_0 = \dfrac{c}{\lambda_0}$已知，则可以得到散射光频率$\nu$的值。$\alpha = 90^\circ$时，反冲电子的动能为：
        \[
        E_\mathrm{k} = E - mc^2 = \sqrt{h^2\nu_0^2 + h^2\nu^2 - 2h^2\nu_0\nu\cos\alpha + m^2c^4} - mc^2
        \]
        故对于$\lambda_0 = 0.100\,\mathrm{nm}$的X射线束，反冲电子的动能为：
        \[
        E_\mathrm{k} = \sqrt{h^2\nu_0^2 + h^2\nu^2 - 2h^2\nu_0\nu\cos\alpha + m^2c^4} - mc^2 = 293.70\,\mathrm{eV}
        \]
        对于$\lambda_0 = 0.0188\,\mathrm{nm}$的$\gamma$射线束，反冲电子的动能为：
        \[
        E_\mathrm{k} = \sqrt{h^2\nu_0^2 + h^2\nu^2 - 2h^2\nu_0\nu\cos\alpha + m^2c^4} - mc^2 = 7538.42\,\mathrm{eV}
        \]

        \item[1.17(3)]
        入射光在碰撞时失去的能量占总能量的比例$k$由下式计算：
        \[
        k = \frac{h\nu_0 - h\nu}{h\nu_0} \times 100\%
        \]
        故对于$\lambda_0 = 0.100\,\mathrm{nm}$的X射线束，入射光在碰撞时失去的能量占总能量的比例为：
        \[
        k = \frac{h\nu_0 - h\nu}{h\nu_0} \times 100\% = 2.369\%
        \]
        对于$\lambda_0 = 0.0188\,\mathrm{nm}$的$\gamma$射线束，入射光在碰撞时失去的能量占总能量的比例为：
        \[
        k = \frac{h\nu_0 - h\nu}{h\nu_0} \times 100\% = 11.43\%
        \]

        \subsection*{补充题}

        \item[1.]
        由Heisenberg不确定性：
        \[
        \Delta x \Delta p \sim \dfrac{h}{4\pi}
        \]
        假设量子效应显著，$\Delta x \sim 1\,\mathrm{m}$，$\Delta p \sim 1\,\mathrm{kg\cdot m/s}$，则$h \sim 4\pi\,\mathrm{J\cdot s}$，即约$10\,\mathrm{J\cdot s}$的数量级。

    \end{itemize}

\end{document}
